\documentclass[11pt, a4paper]{article}
\usepackage[a4paper, margin=2cm]{geometry}
\usepackage{fontspec}
\usepackage[english, russian, bidi=basic, provide=*]{babel}
\babelprovide[import, onchar=ids fonts]{russian}
\babelfont{rm}{Noto Sans}
\usepackage{amsmath, amssymb, amsthm}
\usepackage{booktabs}
\usepackage{enumitem}
\usepackage{hyperref}

\title{\textbf{Полные решения заданий}\\ \large Задание 1. Векторы}
\author{Линейная алгебра}
\date{2024}

\begin{document}

\maketitle

% ==========================================
\section{Раздел 1. Линейные комбинации и уравнения}
% ==========================================

\subsection*{Задание 1.1}
\textbf{Условие:} Найдите линейную комбинацию $7a+3b-4c$ векторов:
$$ a = (3, -1, 3), \quad b = (5, 2, 7), \quad c = (-3, -2, 6) $$

\textbf{Решение:}
Умножим каждый вектор на соответствующий скаляр и выполним покомпонентное сложение:
\begin{align*}
7a &= 7 \cdot (3, -1, 3) = (21, -7, 21), \\
3b &= 3 \cdot (5, 2, 7) = (15, 6, 21), \\
-4c &= -4 \cdot (-3, -2, 6) = (12, 8, -24).
\end{align*}
Складываем полученные векторы:
\begin{align*}
7a + 3b - 4c &= (21 + 15 + 12, \; -7 + 6 + 8, \; 21 + 21 - 24) \\
&= (48, 7, 18).
\end{align*}

\textbf{Ответ:} $(48, 7, 18)$.

---

\subsection*{Задание 1.2}
\textbf{Условие:} Найдите вектор $x$ из следующих уравнений:

\medskip
\textbf{a)} $a + 9b + x - c = 0$, где $a = (3, 5, 1, 6)$, $b = (2, 9, 4, 2)$, $c = (3, 2, 8, 4)$.

\textbf{Решение:}
Выразим вектор $x$:
$$ x = c - a - 9b. $$
Вычислим $9b$:
$$ 9b = 9 \cdot (2, 9, 4, 2) = (18, 81, 36, 18). $$
Теперь посчитаем компоненты вектора $x$:
\begin{align*}
x_1 &= 3 - 3 - 18 = -18, \\
x_2 &= 2 - 5 - 81 = -84, \\
x_3 &= 8 - 1 - 36 = -29, \\
x_4 &= 4 - 6 - 18 = -20.
\end{align*}

\textbf{Ответ:} $x = (-18, -84, -29, -20)$.

\medskip
\textbf{b)} $3.5(a - x) - 2(b + x) + 1.5(c + 3x) = 0$, где $a = (3, 5, 4, 2)$, $b = (2, 3, 4, 2)$, $c = (5, 2, 9, 3)$.

\textbf{Решение:}
Раскроем скобки и сгруппируем слагаемые с неизвестной $x$:
$$ 3.5a - 3.5x - 2b - 2x + 1.5c + 4.5x = 0 $$
$$ (-3.5 - 2 + 4.5)x + 3.5a - 2b + 1.5c = 0 $$
$$ -x + 3.5a - 2b + 1.5c = 0 \implies x = 3.5a - 2b + 1.5c. $$

Вычисляем координаты $x$:
\begin{align*}
x_1 &= 3.5(3) - 2(2) + 1.5(5) = 10.5 - 4 + 7.5 = 14, \\
x_2 &= 3.5(5) - 2(3) + 1.5(2) = 17.5 - 6 + 3 = 14.5, \\
x_3 &= 3.5(4) - 2(4) + 1.5(9) = 14 - 8 + 13.5 = 19.5, \\
x_4 &= 3.5(2) - 2(2) + 1.5(3) = 7 - 4 + 4.5 = 7.5.
\end{align*}

\textbf{Ответ:} $x = (14, 14.5, 19.5, 7.5)$.

---

\subsection*{Задание 1.3}
\textbf{Условие:} Даны векторы $a=(1,2)$, $b=(4,3)$. В результате их линейной комбинации был получен вектор $c=(57,45)$. Найдите коэффициенты.

\textbf{Решение:}
Составим уравнение $\alpha a + \beta b = c$, что приводит к системе линейных уравнений:
$$ \begin{cases} \alpha + 4\beta = 57 \\ 2\alpha + 3\beta = 45 \end{cases} $$
Умножим первое уравнение на $2$:
$$ 2\alpha + 8\beta = 114 $$
Вычтем из него второе уравнение:
$$ (2\alpha + 8\beta) - (2\alpha + 3\beta) = 114 - 45 \implies 5\beta = 69 \implies \beta = \frac{69}{5} = 13.8. $$
Найдём $\alpha$ из первого уравнения:
$$ \alpha = 57 - 4(13.8) = 57 - 55.2 = 1.8 = \frac{9}{5}. $$

\textbf{Ответ:} Коэффициенты: $\alpha = 1.8$ (или $\frac{9}{5}$), $\beta = 13.8$ (или $\frac{69}{5}$).

---

\subsection*{Задание 1.4}
\textbf{Условие:} Найдите значение $\alpha$ такое, что векторы $a=(3,\alpha,2)$ и $b=(2-\alpha,-1,5)$ будут ортогональны (перпендикулярны).

\textbf{Решение:}
Векторы ортогональны тогда и только тогда, когда их скалярное произведение равно нулю:
$$ (a, b) = 0 $$
\begin{align*}
3(2 - \alpha) + \alpha(-1) + 2(5) &= 0 \\
6 - 3\alpha - \alpha + 10 &= 0 \\
16 - 4\alpha &= 0 \implies \alpha = 4.
\end{align*}

\textbf{Ответ:} $\alpha = 4$.

\newpage

% ==========================================
\section{Раздел 2. Нормы, углы и свойства векторов}
% ==========================================

\subsection*{Задание 2.1}
\textbf{Условие:} Посчитайте манхэттенскую ($\ell_1$) и евклидову ($\ell_2$) нормы следующих векторов:
\begin{enumerate}[label=\alph*)]
\item $a = (3, 9, 2, 3, 4, 0, -1, -11, 3)$
\item $b = (12, 13, 23, -22, -34, 31, 21)$
\end{enumerate}

\textbf{Решение:}

\textbf{a)}
\begin{itemize}
\item \textbf{Манхэттенская норма ($\ell_1$):}
$$ \|a\|_1 = |3| + |9| + |2| + |3| + |4| + |0| + |-1| + |-11| + |3| = 3+9+2+3+4+0+1+11+3 = 36. $$
\item \textbf{Евклидова норма ($\ell_2$):}
\begin{align*}
\|a\|_2 &= \sqrt{3^2 + 9^2 + 2^2 + 3^2 + 4^2 + 0^2 + (-1)^2 + (-11)^2 + 3^2} \\
&= \sqrt{9 + 81 + 4 + 9 + 16 + 0 + 1 + 121 + 9} = \sqrt{250} = 5\sqrt{10} \approx 15.811.
\end{align*}
\end{itemize}

\textbf{b)}
\begin{itemize}
\item \textbf{Манхэттенская норма ($\ell_1$):}
$$ \|b\|_1 = |12| + |13| + |23| + |-22| + |-34| + |31| + |21| = 12+13+23+22+34+31+21 = 156. $$
\item \textbf{Евклидова норма ($\ell_2$):}
\begin{align*}
\|b\|_2 &= \sqrt{12^2 + 13^2 + 23^2 + (-22)^2 + (-34)^2 + 31^2 + 21^2} \\
&= \sqrt{144 + 169 + 529 + 484 + 1156 + 961 + 441} = \sqrt{3884} \approx 62.322.
\end{align*}
\end{itemize}

---

\subsection*{Задание 2.2}
\textbf{Условие:} Даны векторы $a=(1,2,3)$ и $b=(4,3,5)$. Найдите угол между ними.

\textbf{Решение:}
Косинус угла $\theta$ между векторами выражается формулой:
$$ \cos \theta = \frac{(a, b)}{\|a\|_2 \|b\|_2} $$
1. Скалярное произведение:
$$ (a, b) = 1 \cdot 4 + 2 \cdot 3 + 3 \cdot 5 = 4 + 6 + 15 = 25. $$
2. Длины векторов:
$$ \|a\|_2 = \sqrt{1^2 + 2^2 + 3^2} = \sqrt{14}, $$
$$ \|b\|_2 = \sqrt{4^2 + 3^2 + 5^2} = \sqrt{16 + 9 + 25} = \sqrt{50} = 5\sqrt{2}. $$
3. Косинус и угол:
$$ \cos \theta = \frac{25}{\sqrt{14} \cdot 5\sqrt{2}} = \frac{5}{\sqrt{28}} = \frac{5}{2\sqrt{7}} \approx 0.94491. $$
$$ \theta = \arccos\left(\frac{5}{2\sqrt{7}}\right) \approx 19.11^\circ \text{ (или } \approx 0.3335 \text{ рад)}. $$

\textbf{Ответ:} $\theta = \arccos\left(\frac{5}{2\sqrt{7}}\right) \approx 19.11^\circ$.

---

\subsection*{Задание 2.3}
\textbf{Условие:} Вектор $x$ длины 100 содержит элементы $x_i \ge 0$ (количество пропусков у $i$-го ученика). Используя манхэттенскую норму, напишите формулу для:
\begin{enumerate}[label=\alph*.]
\item суммы количества пропусков у всех учеников;
\item среднего количества пропусков.
\end{enumerate}

\textbf{Решение:}
Поскольку все $x_i \ge 0$, то $|x_i| = x_i$. Следовательно, манхэттенская норма равна суммы всех элементов:
$$ \|x\|_1 = \sum_{i=1}^{100} |x_i| = \sum_{i=1}^{100} x_i. $$
\begin{enumerate}[label=\alph*.]
\item \textbf{Сумма пропусков:} $S = \|x\|_1$.
\item \textbf{Среднее количество пропусков:} $\bar{x} = \frac{\|x\|_1}{100}$.
\end{enumerate}

---

\subsection*{Задание 2.4}
\textbf{Условие:} Покажите, что следующие уравнения выполняются для любых векторов $a, b$ одинаковой размерности:

\medskip
\textbf{a)} $(a+b)^T(a-b) = \|a\|^2 - \|b\|^2$

\textbf{Доказательство:}
Раскроем скобки по свойству линейности транспонирования и скалярного произведения:
$$ (a+b)^T(a-b) = (a^T + b^T)(a - b) = a^T a - a^T b + b^T a - b^T b. $$
Так как $a^T b$ --- это скаляр, то $a^T b = (a^T b)^T = b^T a$. Поэтому слагаемые $-a^T b$ и $+b^T a$ взаимно уничтожаются:
$$ (a+b)^T(a-b) = a^T a - b^T b = \|a\|^2 - \|b\|^2. \quad \square $$

\medskip
\textbf{b)} $\|a+b\|^2 + \|a-b\|^2 = 2(\|a\|^2 + \|b\|^2)$

\textbf{Доказательство:}
Запишем квадраты евклидовых норм через скалярное произведение:
\begin{align*}
\|a+b\|^2 &= (a+b)^T(a+b) = a^T a + 2a^T b + b^T b = \|a\|^2 + 2a^T b + \|b\|^2, \\
\|a-b\|^2 &= (a-b)^T(a-b) = a^T a - 2a^T b + b^T b = \|a\|^2 - 2a^T b + \|b\|^2.
\end{align*}
Складывая эти два равенства:
$$ \|a+b\|^2 + \|a-b\|^2 = (\|a\|^2 + 2a^T b + \|b\|^2) + (\|a\|^2 - 2a^T b + \|b\|^2) = 2\|a\|^2 + 2\|b\|^2 = 2(\|a\|^2 + \|b\|^2). \quad \square $$

\newpage

% ==========================================
\section{Раздел 3. Системы линейных уравнений (СЛУ) и метод Гаусса}
% ==========================================

\subsection*{Задание 3.1}
\textbf{Условие:} Периметр треугольника равен $3\text{ дм} = 30\text{ см}$. Наибольшая из сторон на $4\text{ см}$ больше средней, а утроенная меньшая сторона равна $\text{средняя сторона} + \text{большая сторона} + 2\text{ см}$. Найдите стороны треугольника методом Гаусса.

\textbf{Решение:}
Обозначим стороны: $x$ --- меньшая, $y$ --- средняя, $z$ --- большая.
Составим уравнения:
1) Периметр: $x + y + z = 30$.
2) $z = y + 4 \implies -y + z = 4$.
3) $3x = y + z + 2 \implies 3x - y - z = 2$.

Запишем расширенную матрицу системы и применитель элементарные преобразования:
$$ \left(\begin{array}{ccc|c} 
1 & 1 & 1 & 30 \\ 
0 & -1 & 1 & 4 \\ 
3 & -1 & -1 & 2 
\end{array}\right) $$

\textbf{Шаг 1:} Из третьей строки вычтем первую строку, умноженную на 3 ($R_3 \leftarrow R_3 - 3R_1$):
$$ \left(\begin{array}{ccc|c} 
1 & 1 & 1 & 30 \\ 
0 & -1 & 1 & 4 \\ 
0 & -4 & -4 & -88 
\end{array}\right) $$

\textbf{Шаг 2:} Из третьей строки вычтем вторую строку, умноженную на 4 ($R_3 \leftarrow R_3 - 4R_2$):
$$ \left(\begin{array}{ccc|c} 
1 & 1 & 1 & 30 \\ 
0 & -1 & 1 & 4 \\ 
0 & 0 & -8 & -104 
\end{array}\right) $$

\textbf{Обратный ход:}
\begin{itemize}
\item Из 3-й строки: $-8z = -104 \implies z = 13\text{ см}$.
\item Из 2-й строки: $-y + 13 = 4 \implies y = 9\text{ см}$.
\item Из 1-й строки: $x + 9 + 13 = 30 \implies x = 8\text{ см}$.
\end{itemize}

\textbf{Ответ:} Стороны треугольника равны $8\text{ см}$, $9\text{ см}$ и $13\text{ см}$.

---

\subsection*{Задание 3.2}
\textbf{Условие:} Решите СЛУ с помощью метода Гаусса:

\medskip
\textbf{a)}
$$ \begin{cases} 3x_1 + 2x_2 + x_3 = 0 \\ 7x_1 + 6x_2 + 5x_3 = 0 \\ 5x_1 + 4x_2 + 3x_3 = 0 \end{cases} $$

\textbf{Решение:}
Расширенная матрица системы:
$$ \begin{pmatrix} 3 & 2 & 1 \\ 7 & 6 & 5 \\ 5 & 4 & 3 \end{pmatrix} $$
\begin{itemize}
\item $R_2 \leftarrow 3R_2 - 7R_1$: получаем строку $(0, 4, 8)$.
\item $R_3 \leftarrow 3R_3 - 5R_1$: получаем строку $(0, 2, 4)$.
\end{itemize}
$$ \begin{pmatrix} 3 & 2 & 1 \\ 0 & 4 & 8 \\ 0 & 2 & 4 \end{pmatrix} \xrightarrow{R_3 \leftarrow 2R_3 - R_2} \begin{pmatrix} 3 & 2 & 1 \\ 0 & 4 & 8 \\ 0 & 0 & 0 \end{pmatrix} $$

Система сведена к ступенчатому виду. Переменная $x_3 = t$ --- свободная.
\begin{itemize}
\item Из 2-й строки: $4x_2 + 8t = 0 \implies x_2 = -2t$.
\item Из 1-й строки: $3x_1 + 2(-2t) + t = 0 \implies 3x_1 - 3t = 0 \implies x_1 = t$.
\end{itemize}

\textbf{Ответ:} $(x_1, x_2, x_3) = (t, -2t, t)$, где $t \in \mathbb{R}$.

\medskip
\textbf{b)}
$$ \begin{cases} -3x_1 + x_2 + x_3 = 0 \\ 5x_1 + x_2 - 2x_3 = 2 \\ -2x_1 - 2x_2 + x_3 = -3 \end{cases} $$

\textbf{Решение:}
$$ \left(\begin{array}{ccc|c} -3 & 1 & 1 & 0 \\ 5 & 1 & -2 & 2 \\ -2 & -2 & 1 & -3 \end{array}\right) $$
\begin{itemize}
\item $R_2 \leftarrow 3R_2 + 5R_1$: $\implies (0, 8, -1 \mid 6)$
\item $R_3 \leftarrow 3R_3 - 2R_1$: $\implies (0, -8, 1 \mid -9)$
\end{itemize}
$$ \left(\begin{array}{ccc|c} -3 & 1 & 1 & 0 \\ 0 & 8 & -1 & 6 \\ 0 & -8 & 1 & -9 \end{array}\right) \xrightarrow{R_3 \leftarrow R_3 + R_2} \left(\begin{array}{ccc|c} -3 & 1 & 1 & 0 \\ 0 & 8 & -1 & 6 \\ 0 & 0 & 0 & -3 \end{array}\right) $$
Третья строка дает несовместное равенство $0 = -3$.

\textbf{Ответ:} Система не имеет решений.

\medskip
\textbf{c)}
$$ \begin{cases} 7x_1 - 5x_2 - 2x_3 - 4x_4 = 8 \\ -3x_1 + 2x_2 + x_3 + 2x_4 = -3 \\ 2x_1 - x_2 - x_3 - 2x_4 = 1 \\ -x_1 + x_3 + 2x_4 = 1 \end{cases} $$

\textbf{Решение:}
Переставим 4-ю строку на 1-е место для удобства:
$$ \left(\begin{array}{cccc|c} -1 & 0 & 1 & 2 & 1 \\ 2 & -1 & -1 & -2 & 1 \\ -3 & 2 & 1 & 2 & -3 \\ 7 & -5 & -2 & -4 & 8 \end{array}\right) $$
Выполним элементарные преобразования над строками:
\begin{itemize}
\item $R_2 \leftarrow R_2 + 2R_1 \implies (0, -1, 1, 2 \mid 3)$
\item $R_3 \leftarrow R_3 - 3R_1 \implies (0, 2, -2, -4 \mid -6)$
\item $R_4 \leftarrow R_4 + 7R_1 \implies (0, -5, 5, 10 \mid 15)$
\end{itemize}
Заметим, что $R_3 = -2R_2$, а $R_4 = 5R_2$. Они полностью обнуляются:
$$ \left(\begin{array}{cccc|c} -1 & 0 & 1 & 2 & 1 \\ 0 & -1 & 1 & 2 & 3 \\ 0 & 0 & 0 & 0 & 0 \\ 0 & 0 & 0 & 0 & 0 \end{array}\right) $$
Свободные переменные: $x_3 = s$, $x_4 = t$.
\begin{itemize}
\item Из 2-й строки: $-x_2 + s + 2t = 3 \implies x_2 = s + 2t - 3$.
\item Из 1-й строки: $-x_1 + s + 2t = 1 \implies x_1 = s + 2t - 1$.
\end{itemize}

\textbf{Ответ:} $(x_1, x_2, x_3, x_4) = (s + 2t - 1, \; s + 2t - 3, \; s, \; t)$, где $s, t \in \mathbb{R}$.

\newpage

% ==========================================
\section{Раздел 4. Бонусные задания и Теория}
% ==========================================

\subsection*{Бонус 1}
\textbf{Условие:} Заданы векторы $a=(1,5,3)$, $b=(6,-4,-2)$, $c=(0,-5,7)$, $d=(-20,27,35)$. Подобрать числа $\alpha, \beta, \gamma$, чтобы векторы $\alpha a, \beta b, \gamma c, d$ образовывали замкнутую ломаную линию.

\textbf{Решение:}
Условие замкнутости ломаной означает, что сумма векторов равна нулевому вектору:
$$ \alpha a + \beta b + \gamma c + d = 0 \implies \alpha a + \beta b + \gamma c = -d. $$
Подставляем координаты:
$$ \alpha(1, 5, 3) + \beta(6, -4, -2) + \gamma(0, -5, 7) = (20, -27, -35). $$
Составим систему уравнений:
$$ \begin{cases} \alpha + 6\beta = 20 \\ 5\alpha - 4\beta - 5\gamma = -27 \\ 3\alpha - 2\beta + 7\gamma = -35 \end{cases} $$

Из 1-го уравнения выразим $\alpha = 20 - 6\beta$ и подставим во 2-е и 3-е:
$$ 5(20 - 6\beta) - 4\beta - 5\gamma = -27 \implies 100 - 34\beta - 5\gamma = -27 \implies 34\beta + 5\gamma = 127 $$
$$ 3(20 - 6\beta) - 2\beta + 7\gamma = -35 \implies 60 - 20\beta + 7\gamma = -35 \implies 20\beta - 7\gamma = 95 $$

Умножим первое уравнение на 7, а второе на 5:
$$ \begin{cases} 238\beta + 35\gamma = 889 \\ 100\beta - 35\gamma = 475 \end{cases} $$
Сложим эти уравнения:
$$ 338\beta = 1364 \implies \beta = 4. $$
Найдём $\gamma$: $20(4) - 7\gamma = 95 \implies 80 - 7\gamma = 95 \implies -7\gamma = 15 \implies \gamma = -2$.
Найдём $\alpha$: $\alpha = 20 - 6(4) = -4$.

\textbf{Ответ:} $\alpha = -4$, $\beta = 4$, $\gamma = -2$.

---

\subsection*{Бонус 2}
\textbf{Условие:} Найдите решение системы $Ax = 12x$, где
$$ A = \begin{pmatrix} 6 & 4 & 3 \\ 6 & 0 & 9 \\ 0 & 8 & 0 \end{pmatrix} \quad \text{и} \quad \sum_{i=1}^3 x_i = 1. $$

\textbf{Решение:}
Уравнение $Ax = 12x$ перепишется в виде однородной системы $(A - 12I)x = 0$:
$$ \begin{pmatrix} -6 & 4 & 3 \\ 6 & -12 & 9 \\ 0 & 8 & -12 \end{pmatrix} \begin{pmatrix} x_1 \\ x_2 \\ x_3 \end{pmatrix} = \begin{pmatrix} 0 \\ 0 \\ 0 \end{pmatrix} $$
Приведём матрицу системы к ступенчатому виду:
\begin{itemize}
\item Прибавим 1-ю строку ко 2-й ($R_2 \leftarrow R_2 + R_1$):
$$ \begin{pmatrix} -6 & 4 & 3 \\ 0 & -8 & 12 \\ 0 & 8 & -12 \end{pmatrix} $$
\item Прибавим 2-ю строку к 3-й ($R_3 \leftarrow R_3 + R_2$):
$$ \begin{pmatrix} -6 & 4 & 3 \\ 0 & -8 & 12 \\ 0 & 0 & 0 \end{pmatrix} $$
\end{itemize}
Из 2-й строки: $-8x_2 + 12x_3 = 0 \implies x_2 = \frac{12}{8} x_3 = 1.5 x_3$. \\
Из 1-й строки: $-6x_1 + 4(1.5x_3) + 3x_3 = 0 \implies -6x_1 + 9x_3 = 0 \implies x_1 = 1.5 x_3$.

Используем условие нормировки $x_1 + x_2 + x_3 = 1$:
$$ 1.5x_3 + 1.5x_3 + x_3 = 1 \implies 4x_3 = 1 \implies x_3 = \frac{1}{4} = 0.25. $$
Тогда $x_1 = 1.5 \cdot 0.25 = \frac{3}{8} = 0.375$ и $x_2 = \frac{3}{8} = 0.375$.

\textbf{Ответ:} $x = \left(\frac{3}{8}, \frac{3}{8}, \frac{1}{4}\right)^T = (0.375, 0.375, 0.25)^T$.

---

\subsection*{Теория}

\textbf{a) В каком случае векторы будут ортогональны (перпендикулярны)?}
\begin{proof}[Ответ]
Два вектора $a$ и $b$ ортогональны тогда и только тогда, когда их скалярное произведение равно нулю: $(a, b) = 0$ (или $a^T b = 0$).
\end{proof}

\textbf{b) Приведите свой пример того, что можно представить в виде вектора.}
\begin{proof}[Ответ]
В качестве примера можно привести профиль пользователя в рекомендательной системе (например, фильмам или товарам), где каждый компонент вектора отвечает за степень интереса к определенной категории: $x = (\text{Фантастика}, \text{Комедия}, \text{Драма}) = (0.9, 0.1, 0.4)$.
\end{proof}

\textbf{c) Если в СЛУ одна переменная во всех уравнениях имеет нулевой коэффициент, то сколько решений может иметь вся СЛУ и почему?}
\begin{proof}[Ответ]
Вся СЛУ может иметь либо \textbf{0 решений} (быть несовместной), либо \textbf{бесконечно много решений}. 

\textit{Объяснение:} Если переменная $x_k$ входит во все уравнения с нулевым коэффициентом, её значение никак не влияет на выполнение ни одного из уравнений системы.
1. Если подсистема из остальных переменных не имеет решений (содержит противоречие), то вся СЛУ не имеет решений ($0$ решений).
2. Если подсистема совместна и имеет хотя бы одно решение для остальных переменных, то $x_k$ может принимать абсолютно любое значение из поля скаляров ($\mathbb{R}$), что даёт бесконечное множество решений. В силу того, что $x_k$ гарантированно является свободной переменной, ровно 1 единственное решение такая СЛУ иметь не может ни при каких условиях.
\end{proof}

\end{document}